% ======================================================================
% Development Report for the compass_sim Project — V16
% Generated on 05/07/2026. Covered history: 15/06/2026 to 05/07/2026.
% Compile with: pdflatex relatorio_projeto_V16_EN.tex (twice for references)
% ======================================================================
\documentclass[11pt,a4paper]{article}

\usepackage[utf8]{inputenc}
\usepackage[T1]{fontenc}
\usepackage[english]{babel}
\usepackage{geometry}
\geometry{margin=2.5cm}
\usepackage{graphicx}
\usepackage{booktabs}
\usepackage{longtable}
\usepackage{amsmath}
\usepackage{amssymb}
\usepackage{siunitx}
\usepackage{hyperref}
\usepackage{xcolor}
\hypersetup{colorlinks=true, linkcolor=blue!50!black, urlcolor=blue!50!black,
            citecolor=blue!50!black}
\usepackage{caption}
\captionsetup{font=small, labelfont=bf}

\newcommand{\code}[1]{\texttt{#1}}

\title{\textbf{Simulation of dipolar magnetic-needle arrays:\\
development report for the \code{compass\_sim} project}\\[1ex]
\large Report version V16}
\author{J.~P.~Sinnecker}
\date{July 5, 2026\\[0.5ex]
\small Covered period: 15/06/2026 -- 05/07/2026}

\begin{document}
\maketitle

\begin{abstract}
\noindent
This report documents the full development of the \code{compass\_sim} project: a Python simulator for two-dimensional arrays of magnetic compass needles interacting through complete dipolar coupling, with inertial Newtonian dynamics in SI units. The project began as a pedagogical demonstration of the analogy between compass arrays and magnetic materials, including hysteresis and domain formation, and evolved over six development sessions into a research instrument with publication potential. The report records the project decisions, problems encountered and solutions adopted, the code-version history (V20--V50), the research and publication strategy built around a gap identified in the literature, and the results of the pilot campaign on ``damping~$\times$~hysteresis/avalanches''. The pilot already shows the central signal sought: a systematic increase of the avalanche exponent $\tau$ with the quality factor $Q$.
\end{abstract}

\tableofcontents
\newpage

% ======================================================================
\section{Project objective}
% ======================================================================

The project now has two layered objectives, reflecting its historical evolution.

\paragraph{Original objective (pedagogical/demonstrative).}
To build a visual simulation of an array of compass needles that rotate only under the magnetic field generated by each needle on its neighbors, seeking the collective directional equilibrium. The simulation provides a direct and intuitive mechanical analogy with real magnetic materials, capable of displaying hysteretic behavior and magnetic-domain formation from a top view and exportable as video.

\paragraph{Current objective (scientific).}
To use the simulator, which implements real Newtonian \emph{inertial} dynamics, a rare feature in this literature, to address a nonequilibrium statistical-physics question with publication potential:

\begin{quote}\itshape
Does the damping coefficient control the universality class of avalanches and the shape of the hysteresis loop in classical dipolar arrays with long-range and anisotropic coupling? Does this response depend on the lattice geometry (square, triangular, honeycomb), that is, on geometric frustration?
\end{quote}

The literature has established that underdamped and overdamped systems belong to different universality classes in short-range or mean-field models, such as inertial Kuramoto models, damped Ising models, and sheared amorphous solids~\cite{dahmen2011}. The identified gap is that no one has tested whether this picture survives real dipolar coupling, which is long-ranged ($\sim 1/r^3$) and anisotropic, in two-dimensional arrays with continuous angular dynamics and physical inertia. As far as the literature review reached, the simulator developed here is the only instrument built specifically for this question.

% ======================================================================
\section{The physical model}
% ======================================================================

Each needle $i$ is represented as a point magnetic dipole of moment $m$ fixed at a lattice position \vec r_i, free to rotate in the plane around a central pin with no contact friction, but with viscous damping from air. The dynamics is Newton's second law for rotation:
\begin{equation}
I\,\ddot\theta_i \;=\; \tau_i^{\mathrm{dip}} + \tau_i^{\mathrm{ext}}
- b\,\dot\theta_i,
\label{eq:eom}
\end{equation}
where $I$ is the needle's moment of inertia, $b$ is the viscous damping coefficient, and the torques come from $\vec\tau = \vec m \times \vec B$. The dipolar field acting on needle $i$ sums the contributions from \emph{all} other needles, using complete coupling with computational cost $\mathcal{O}(K^2)$ for $K$ needles:
\begin{equation}
\vec B_j(\vec r_i) = \frac{\mu_0}{4\pi}\,
\frac{3(\vec m_j \cdot \hat r)\hat r - \vec m_j}{r^3},
\qquad \vec r = \vec r_i - \vec r_j .
\end{equation}

The integration uses a Velocity--Verlet scheme with time step $\Delta t$ proportional to the natural period $T_0 = 2\pi/\omega_0$, where
\begin{equation}
\omega_0 = \sqrt{\frac{m\,B_{\mathrm{ref}}}{I}},
\qquad
B_{\mathrm{ref}} = \frac{\mu_0}{4\pi}\frac{2m}{r_{nn}^{3}}
\end{equation}
is the small-oscillation frequency in the reference dipolar field between nearest neighbors, separated by $r_{nn}$. The dynamical regime is characterized by the quality factor
\begin{equation}
Q = \frac{\omega_0 I}{b},
\end{equation}
with $Q \gg 1$ underdamped (oscillatory) and $Q \ll 1$ overdamped (monotonic relaxation). This is the central control parameter of the scientific study.

\paragraph{Real physical parameters.}
A structuring decision was to use SI units and derive the parameters from the real needle geometry rather than arbitrary constants:
\begin{itemize}
\item \textbf{Moment of inertia}: computed for a \SI{0.4}{mm}-thick steel sheet shaped as a diamond needle, for example $I \approx \SI{3.8e-11}{kg.m^2}$ and mass $\approx$ \SI{0.017}{g} for a \SI{0.5}{cm} needle;
\item \textbf{Magnetic moment}: assuming steel saturated along the long axis, $m = M_s V$, with $M_s = \SI{1.59e6}{A/m}$, corresponding to ordinary carbon steel with $B_{\mathrm{sat}} \approx \SI{2.0}{T}$;
\item \textbf{Lattice}: parameterized by the radius $R$ of the circle circumscribing each needle ($r_{nn} = 2R$ in all geometries) and by the fraction \code{needle\_frac} of the diameter occupied by the needle.
\end{itemize}
Both automatic calculations can be overridden by explicit user parameters.

\paragraph{Simulator capabilities (current state, V50).}
The simulator supports three lattice geometries (square, triangular, honeycomb); \textbf{six} external-field modes: static; pulse with explicit timing (delay $\Delta t$ $\to$ field for $t_{\mathrm{pulse}}$ $\to$ zero field until equilibrium, with the legacy saturation/plateau criterion if $t_{\mathrm{pulse}}$ is not given); five-segment hysteresis ($0 \to +B \to 0 \to -B \to 0 \to +B$) with linear or log-symmetric spacing; sinusoidal field; and step modes \code{step\_pos} ($B{=}0$ for $\Delta t$, then field kept on) and \code{step\_neg} (field from $t{=}0$, removed at $\Delta t$). It also supports periodic boundary conditions with lattice images; real-time monitoring of the alignment order parameter $S \in [0,1]$; a real-equilibrium criterion (small angular velocity \emph{and} small mean torque) for pulse/step modes; optional GPU backend through CuPy; MP4 video export, summary images, and CSV time series $(t, B, M, S)$; magnetic-domain visualization through halo coloring; a precomputed dipolar tensor (V40); and performance instrumentation.

% ======================================================================
\section{Development chronology}
% ======================================================================

Development took place over six intensive sessions (15--18/06/2026), followed by successive consolidation sessions (04--05/07/2026). This section summarizes the decisions and results from each phase; Table~\ref{tab:versoes} consolidates the version record.

\subsection{Phase 1 --- simulation core (15/06)}

The first session established the core: traditional compass needles shaped as diamonds (blue north half, white south half), viewed from above; fully dipolar interaction among all needles; and the search for directional equilibrium. Successive requests added lattice parameters (spacing, $N \times M$ count, square/triangular choice), honeycomb geometry (hexagon vertices), a uniform external field with controllable magnitude and direction, comments documenting all code sections, saving in the current directory, and the video pipeline (PNG frames at regular intervals assembled into MP4 through ffmpeg).

Two important physical decisions were made in this phase: (i) \textbf{SI units} for field and magnetic moment, abandoning arbitrary units; and (ii) \textbf{inertial dynamics}: negligible pin contact friction, but real moment of inertia, so the needle \emph{oscillates} around equilibrium before stabilizing, and this must be visible in the animation. This second decision, originally adopted for visual realism, later became the scientific differential of the project, since inertial dynamics is precisely what is missing in the literature on dipolar arrays.

\subsection{Phase 2 --- control, physical time, and honeycomb geometry (16--17/06)}

This session was dominated by three fronts.

\textbf{(a) Interface and control}: visual indication of external-field direction and intensity in the image; automatic incremental numbering of videos (\code{name0000.mp4}, \code{0001}, \dots); controlled interruption through \mbox{Ctrl+I} or a physical criterion ($S = 1.00$); and explanation/exposure of the damping parameter and the quality factor $Q$.

\textbf{(b) Time semantics}: \code{t\_sim} was redefined to mean the simulated \emph{real physical time}, namely the sum of the integrated $\Delta t$ values, equivalent to what a stopwatch would measure while observing the real system, rather than the wall time of the execution. This distinction led to successive corrections until the behavior became consistent.

\textbf{(c) Honeycomb correction}: the initial implementation produced distorted networks that were not truly hexagonal. After several attempts, the adopted solution was to construct the honeycomb as a triangular lattice with systematic removal of the central site of each hexagon, and to parameterize all geometries by the radius $R$ of the circle circumscribing the needle. This guarantees a uniform $r_{nn} = 2R$ across geometries, which later proved essential for clean scientific comparison among them, with the same nominal $Q$ in all three lattices.

\subsection{Phase 3 --- dynamic-field modes and data export (17/06)}

This phase transformed the demonstrator into a measurement instrument:

\begin{itemize}
\item \textbf{Hysteresis mode}: field linearly swept through the full cycle. An important protocol correction was requested and implemented: the correct cycle has \emph{five} segments ($0 \to +B_{\max} \to 0 \to -B_{\max} \to 0 \to +B_{\max}$), not three;
\item \textbf{Sinusoidal mode}: $B(t) = B_{\max}\sin(2\pi f t)$ in the chosen direction;
\item \textbf{Pulse mode}: field applied until $S > 0.99$ \emph{or} an $S$ plateau is reached (variation $<5\%$ over 10 iterations), then set to zero to observe free relaxation. The plateau criterion was added after the saturation-only criterion proved insufficient;
\item \textbf{CSV export}: $(t, B, M, S)$ series with magnetization projected along the field direction, using the same numbering scheme as the videos. This became the backbone of the subsequent scientific analysis;
\item Robustness fixes: degenerate case $N = M = 1$, terminal output with corrupted line feeds, and image resolution for large arrays ($\geq 20 \times 20$).
\end{itemize}

Periodic boundary conditions (PBC) were introduced at the end of this phase as an on/off parameter, then completed in the next phase with image replicas in the dipolar-field calculation (default: one copy in each direction).

\subsection{Phase 4 --- version control, GPU, and automatic parameters (17/06)}

An organizational milestone was reached: by request, all previous versions were retroactively consolidated (V1--V19), and a \textbf{formal version record} was established from \textbf{V20} onward, with an increment for each change. This practice has been maintained since then and also adopted for this report.

This phase included refinement of the pulse mode; repositioning of on-screen indicators (field at the top right, with a dedicated upper margin to avoid overlap with needles); and \textbf{GPU support} via CuPy, motivated by the availability of an NVIDIA RTX~3090 on the computing server. The investigation of why GPU utilization remained low led, in the next phase, to complete vectorization of the torque loop, eliminating interpreted nested Python loops with $\mathcal{O}(K^2)$ cost, and to the honest conclusion that for small and medium arrays the dominant cost was not computation but rendering. This redirected the optimization effort.

\subsection{Phase 5 --- performance, interruption, and magnetic domains (17--18/06)}

\begin{itemize}
\item \textbf{Performance instrumentation}: report of integration time, steps per second, and ms/step at the end of each run; progress bar with percentage; explicit \code{--gpu} flag;
\item \textbf{Automatic moment of inertia} from the real geometry (\SI{0.4}{mm} steel sheet), replaceable by parameter;
\item \textbf{Interruption handling}: \mbox{Ctrl+I} interrupts gracefully (saving video/data); \mbox{Ctrl+C} aborts immediately. The investigation revealed that the default handler was masked by the integration loop, and this was corrected (exit code 130);
\item \textbf{Rendering optimization (V35)}: replacement of thousands of individual matplotlib calls by vectorized collections (\emph{PatchCollections} + \emph{quiver}), with a measured speedup of \textbf{17--30$\times$}: a $25\times25$ array that exceeded 30~s of rendering dropped to 3.9~s, and $40\times40$ (previously impractical) dropped to 14~s, with bitwise validation of polygon geometry and physics CSV;
\item \textbf{Automatic magnetic moment (V37)} from the saturated volume ($M_s = \SI{1.59e6}{A/m}$, researched for ordinary carbon steel), and correction of a regression where \code{--progress\_bar 0} also suppressed periodic status messages;
\item \textbf{Magnetic-domain visualization (V38)}: halo coloring by domain. Neighboring needles whose directions differ by less than an angular tolerance (\code{--domain\_tol}, default $15^\circ$) are grouped through \emph{union--find}; small domains ($<\max(3;\,2\%)$ of the needles) are painted neutral gray so the main domains stand out (Figure~\ref{fig:dominios}). This feature began as visualization and became a measurement instrument: the same routine feeds the domain-size statistics of the scientific campaign.
\end{itemize}

\begin{figure}[t]
\centering
\includegraphics[width=0.72\textwidth]{figs/halo_domains_v3.png}
\caption{Magnetic-domain visualization (V38): square $24\times24$ array after long relaxation. The five main domains receive distinct colors and small domains appear in gray. The same labeling routine feeds the domain statistics of the scientific campaign.}
\label{fig:dominios}
\end{figure}

\subsection{Phase 6 --- from tool to research (18/06 and 04--05/07)}

With the instrument mature, the focus shifted to science:

\begin{enumerate}
\item \textbf{Literature review}, detailed in Section~\ref{sec:pesquisa}, identifying the gap and study design;
\item \textbf{Publication strategy}, including ACS journals, the CAPES agreement, and high-impact alternatives;
\item \textbf{\code{--make\_images} (V39)}: it was found that, even without \code{--video}, the program unconditionally generated 4--5 summary PNG images per run. This cost is negligible in large arrays, but dominant in small arrays used for sweeps. The new parameter suppresses all image generation while preserving data CSVs, with a measured speedup of $\sim 3\times$ in small runs (2.03~s $\to$ 0.65~s for an $8\times8$ array);
\item \textbf{Campaign infrastructure}: the scripts \code{damping\_sweep.py} (orchestration of damping~$\times$~seed~$\times$~geometry sweeps, calling the simulator programmatically) and \code{damping\_sweep\_analysis.py} (full statistical analysis), both validated end to end (Section~\ref{sec:campanha});
\item \textbf{Performance optimizations implemented (04/07)}: the precomputed dipolar tensor (\textbf{V40}) and campaign parallelization with incremental \emph{manifest} and resume support (\code{damping\_sweep} V02). Details and measurements are in Section~\ref{sec:melhorias}. The measured V40 speedup is \textbf{35--45$\times$} per run, reducing the full campaign from $\sim$7--8~h to $\sim$15~min on a single core;
\item \textbf{New field protocols and refinements (V41--V45)}: hysteresis with log-symmetric spacing; pulse with explicit timing ($\Delta t$ + $t_{\mathrm{pulse}}$ + torque equilibrium); step modes \code{step\_pos}/\code{step\_neg}; \code{--t\_sim\_full} flag (run to the end without early stopping); halos also in the comparison figure (V43); a two-line live progress panel (V44); and hysteresis video with a dual progress panel (V45).
\item \textbf{Layout and terminal updates (V46--V50)}: docstring alignment (V46); autoscale parameter \code{--hyst\_autoscale} (V47); bugfixes and autoscale options (V48); HUD-needle overlap resolution and signed field bar updates (V49); and progress bar English translation alongside a three-line console print layout (V50).
\end{enumerate}

% ======================================================================
\section{Research and publication strategy}
\label{sec:pesquisa}
% ======================================================================

\subsection{The gap in the literature}

The review covered three fronts. \textbf{(i)}~Classical compass-needle arrays: the founding work (square $22\times22$ lattice, Phys.\ Rev.\ B \textbf{58}, 9238, 1998~\cite{prb1998}) uses equilibrium Monte Carlo with a random field mimicking thermal noise, not Newtonian dynamics. \textbf{(ii)}~Classical dipoles in honeycomb/kagome/triangular lattices (2012--2025): Luttinger--Tisza methods, energy minimization, spin waves, and Monte Carlo with \emph{parallel tempering}, all focused on ground states and thermal \emph{equilibrium} transitions. \textbf{(iii)}~\emph{Artificial spin ice}~\cite{nisoli2013}: real temporal dynamics is studied, but through discrete macrospins (Ising-like) or experimental data (MOKE/MFM), almost never with continuous angular dynamics and explicit mechanical inertia.

The question of how damping controls hysteresis and avalanches is active and publishable in adjacent fields. It is established that underdamped and overdamped regimes form different universality classes in short-range or mean-field systems~\cite{dahmen2011}, but this has not been answered for complete dipolar coupling (long-range, anisotropic) with real spatial geometry. This is the gap occupied by the project.

\subsection{Study design and decision criterion}

The central study sweeps damping over $\sim 2.5$ orders of magnitude in $Q$ (from $Q \approx 15$, underdamped, to $Q \approx 0.05$, overdamped), with multiple initial-noise seeds in the three geometries. It measures hysteresis-loop area, coercive field, remanence, avalanche statistics (power-law exponent $\tau$ in $P(s) \sim s^{-\tau}$, fitted by maximum likelihood using the Clauset--Shalizi--Newman method~\cite{clauset2009}), and domain-size distribution in free relaxation.

The decision criterion was defined \emph{a priori}:
\begin{itemize}
\item \textbf{Result A (incremental, solid)}: $\tau$ approximately constant with $Q$ (only scales/cutoffs change), consistent with known phenomenology and publishable in the standard target journals;
\item \textbf{Result B (surprising)}: $\tau$ varies systematically with $Q$, or there is a qualitative difference between geometries, signaling a breakdown of universality due to anisotropic long-range dipolar coupling. This would justify attempting PRL or a similar venue.
\end{itemize}

\subsection{Target journals and APC funding}

\begin{itemize}
\item \textbf{CAPES--ACS agreement}: it was confirmed that the CAPES transformative \emph{Read \& Publish} agreement covers \textbf{100\% of the APC in any ACS journal} (CC~BY license) for corresponding authors from consortium institutions, with eligibility checked during submission via ORCID, removing cost as a decision factor;
\item \textbf{Primary target}: \emph{ACS Materials Au} (IF~6.5, Q2), the highest-impact fully \emph{open access} ACS option. The paper must be framed as a model for \emph{artificial spin ice} and patterned nanomagnets, which was incorporated into the introduction design;
\item \textbf{Backup target}: \emph{ACS Physical Chemistry Au} (IF~4.3, Q2), whose scope explicitly includes condensed matter and is safer in fit;
\item \textbf{High-impact alternatives}: \emph{Physical Review Letters} is \emph{outside} the CAPES agreement (APS is not a participating publisher); \emph{Nature Communications}, being \emph{gold OA}, is excluded from \emph{Read \& Publish} agreements (APC $\sim$ USD~6,190 separately). Both are justified only for a type-B result;
\item \textbf{Other CAPES publishers} (Springer Nature $\sim$1,750 hybrid titles, Elsevier $\sim$160, Wiley, IEEE, RSP, ACM): a partial review did not find, up to letter~H of the Springer list, a title with better fit than the ACS targets; natural Elsevier candidates (\emph{Physica A}, \emph{JMMM}) await verification in the eligible-title list.
\end{itemize}

% ======================================================================
\section{Campaign infrastructure and pilot study}
\label{sec:campanha}
% ======================================================================

\subsection{The two scripts}

\textbf{\code{damping\_sweep.py}} orchestrates the campaign by calling the simulator programmatically without modifying the source code. It builds the log-spaced damping grid in $Q$ from the real physics (inertia and moment derived from geometry), runs Step~1 (hysteresis) and Step~2 (free relaxation with domain statistics) for each geometry~$\times$~$Q$~$\times$~seed combination, and saves per-run CSVs, JSON metadata (including $Q$, $\omega_0$, $B_{\mathrm{ref}}$, domain distribution), and a consolidated \code{manifest.csv}.

\textbf{\code{damping\_sweep\_analysis.py}} reads \code{manifest.csv} and produces: loop area (closed-line integral), coercivity and remanence (interpolation at zero crossings); avalanche detection in $dM/dB$ within each monotonic ramp segment (robust threshold of $4\times$MAD); power-law fitting by MLE with $x_{\min}$ selected by KS distance~\cite{clauset2009}, \emph{aggregating avalanches from all seeds} for each $(geometry, Q)$ point before fitting, a methodological choice required for reliable statistics; and final plots, including the central $\tau(Q)$ test by geometry.

Validations performed: area against an exact rectangular loop; $B_c$/$M_r$ against a synthetic loop with known values; avalanche detection against injected steps (size recovery with $\sim$1\% error); MLE against a synthetic sample with exact $\tau = 2.5$ (recovered $2.516 \pm 0.021$), through both code paths (the \code{powerlaw} package, with a real bug in the installed version corrected, and the in-house fallback implementation, with identical results).

\subsection{Pilot-study results}

A reduced pilot (array $14\times14 = 196$ needles, 5 $Q$ values, 3 seeds, 3 geometries, shortened cycle) validated the full pipeline and already displays the central signal sought (Figure~\ref{fig:tau} and Table~\ref{tab:piloto}): the exponent $\tau$ \textbf{increases systematically with $Q$} in all three geometries, from $\tau \approx 1.2$--$1.3$ in the \emph{crossover} ($Q \approx 0.87$) to $\tau \approx 1.7$--$1.8$ in the underdamped regime ($Q = 15$). For the two most overdamped $Q$ values, the shortened pilot cycle does not generate enough avalanches for fitting, which is the correct pipeline behavior: it reports unavailability instead of a spurious number.

\begin{table}[h]
\centering
\small
\caption{Pilot: avalanche exponent $\tau$ (MLE, avalanches aggregated across 3 seeds) by geometry and quality factor. ``---'' indicates insufficient statistics ($n_{\mathrm{tail}} <$ minimum) in the shortened pilot cycle.}
\label{tab:piloto}
\begin{tabular}{@{}lccc@{}}
\toprule
$Q$ & square & triangular & honeycomb \\
\midrule
15.0  & $1.81 \pm 0.15$ & $1.69 \pm 0.13$ & $1.73 \pm 0.15$ \\
3.60  & $1.62 \pm 0.10$ & $1.51 \pm 0.08$ & $1.60 \pm 0.12$ \\
0.866 & $1.19 \pm 0.04$ & $1.29 \pm 0.06$ & --- \\
0.208 & --- & --- & --- \\
0.050 & --- & --- & --- \\
\bottomrule
\end{tabular}
\end{table}

\begin{figure}[t]
\centering
\includegraphics[width=0.8\textwidth]{figs/tau_vs_Q.png}
\caption{Central test of the study (pilot): avalanche power-law exponent $\tau$ as a function of the quality factor $Q$, by geometry. The systematic increase of $\tau$ with $Q$, if confirmed in the full campaign with adequate statistics, is the candidate signature that the universality class depends on the damping regime under dipolar coupling.}
\label{fig:tau}
\end{figure}

\begin{figure}[t]
\centering
\includegraphics[width=0.75\textwidth]{figs/area_vs_Q.png}
\caption{Pilot: hysteresis-loop area as a function of $Q$, by geometry. The area increases monotonically when entering the overdamped regime (more dissipation per cycle), with the three geometries qualitatively similar in this metric, unlike domain size, for which the triangular lattice shows non-monotonic behavior to be investigated.}
\label{fig:area}
\end{figure}

Two secondary findings from the pilot are worth recording: (i) loop area and coercive field increase monotonically when entering the overdamped regime (Figure~\ref{fig:area}); (ii) the mean domain size in the \emph{triangular} lattice shows non-monotonicity with $Q$ that is not seen in the other geometries, exactly the kind of frustration~$\times$~inertia interaction identified in the study design as a type-B candidate result. Both require the full campaign for confirmation.

\subsection{Full campaign (ready to run)}

Defined parameters: $30\times30$ lattice (900 needles), 8 log-spaced $Q$ values in $[0.05;\,15]$, 5 seeds, 3 geometries $\Rightarrow$ 120 hysteresis runs + 120 relaxation runs. Cost with V39 code: $\sim$2--2.5~min per run ($\approx 7$--8~h serial campaign). \textbf{With V40 optimization} (precomputed tensor, Section~\ref{sec:melhorias}): $\sim$3.3--3.4~s per run $\Rightarrow$ \textbf{$\sim$15~min serial campaign on a single core}, and only a few minutes with \code{--n\_workers} on the \code{jps@mumax} server, with no GPU required.

% ======================================================================
\section{Step-by-step roadmap for the proposed studies}
\label{sec:roteiro}
% ======================================================================

This section details, in operational format, the four studies identified as publication candidates, from the central study to the satellites. For each study it specifies the question, simulation steps, analysis, expected result, and editorial destination. All use only already implemented capabilities (V39 + campaign scripts), except where indicated.

\subsection{Study 1 (central): damping as a control parameter for hysteresis and avalanches}

\textbf{Question.} Does the quality factor $Q$ control the universality class of avalanches and the shape of the hysteresis loop under complete dipolar coupling? Does the answer depend on geometry?

\begin{enumerate}
\item \textbf{Calibration} ($\sim$30~min): isolated needle in pulse mode; fit the free-decay envelope and validate that the measured $Q$ agrees with $Q = \omega_0 I / b$. This ensures that the damping grid actually spans both regimes;
\item \textbf{Hysteresis campaign}: $30\times30$ array; 8 log-spaced $Q$ values in $[0.05;\,15]$; 5 seeds; 3 geometries; five-ramp cycle with $B_{\max} = 8\,B_{\mathrm{ref}}$ and duration $\propto T_0$ (120 runs, $\sim$4.5~h on serial CPU without V40);
\item \textbf{Relaxation campaign}: same combinations with $B = 0$, starting from the same noise, up to the $S$ plateau; domain statistics in the final state (120 runs, $\sim$3~h without V40);
\item \textbf{Analysis} (\code{damping\_sweep\_analysis.py}): area, $B_c$, and $M_r$ vs.\ $Q$ by geometry; $\tau(Q)$ by MLE with avalanches aggregated across seeds; domain size/number vs.\ $Q$;
\item \textbf{Robustness} (mandatory before conclusions): convergence in $\Delta t$ at the extremes of $Q$; sizes $N = 20$ and $40$ at $Q \approx 1$; comparison with/without PBC;
\item \textbf{Decision and writing}: A/B criterion from Section~\ref{sec:pesquisa}; framing as a ``model for \emph{artificial spin ice} / patterned nanomagnets'' in the introduction.
\end{enumerate}

\textbf{Expected result.} The pilot suggests increasing $\tau$ with $Q$ in the three geometries. If confirmed with full statistics and surviving robustness diagnostics, this is a type-B candidate result. \textbf{Destination:} ACS Materials Au (result A) or an attempt at PRL (result B), with ACS Physical Chemistry Au as backup.

\subsection{Study 2 (satellite): geometric frustration and domain-formation kinetics}

\textbf{Question.} Does geometric frustration (triangular $>$ honeycomb $>$ square) accelerate or slow down domain coarsening? Is the non-monotonicity of domain size with $Q$ seen in the triangular-lattice pilot real?

\begin{enumerate}
\item Long free relaxation ($B = 0$) in the three geometries, same seeds, for 3--4 representative $Q$ values;
\item Record $S(t)$ and label domains at intermediate times. This \emph{requires minor instrumentation}: domain statistics are currently computed only in the final state; periodic labeling \emph{snapshots} must be added along the trajectory;
\item Extract the time to the $S$ plateau and the growth law of the characteristic domain size $L(t) \sim t^{n}$, comparing exponent $n$ across geometries and damping regimes.
\end{enumerate}

\textbf{Destination:} may form part of the central article (domain section) or, if the frustration~$\times$~inertia interaction is rich, a separate article in ACS Physical Chemistry Au. Computational cost is low, since no field ramps are required.

\subsection{Study 3 (satellite): memory of the \emph{director} after a pulse}

\textbf{Question.} After a saturating field pulse, does the \emph{director} (emergent preferred direction, concept from the 1998 work~\cite{prb1998}) formed during relaxation retain memory of the pulse direction, or is it erased and redetermined by the initial noise? Equilibrium Monte Carlo cannot answer this, since it requires a real time trajectory.

\begin{enumerate}
\item Pulse mode with field direction $\varphi$ swept, for example from $0^\circ$ to $90^\circ$ in $15^\circ$ steps, several seeds, and 2--3 $Q$ values;
\item Measure the direction of the final \emph{director} (mean orientation of the main domains) as a function of $\varphi$;
\item Quantify the directional correlation $\langle \cos 2(\theta_{\mathrm{dir}} - \varphi) \rangle$ and its dependence on $Q$. The interesting hypothesis is that memory is stronger in the overdamped regime, without oscillations that ``shake'' the system out of the valley selected by the pulse.
\end{enumerate}

\textbf{Destination:} connects to the recent literature on \emph{clocked dynamics} and training in \emph{artificial spin ice}; a good candidate for a second article. Moderate cost.

\subsection{Study 4 (satellite): boundary conditions and finite size in the domain scale}

\textbf{Question.} How much does the free boundary distort domain and avalanche statistics relative to the periodic limit? Long-range dipolar systems are notoriously sensitive to sample shape; the answer is methodologically useful for the whole community simulating these arrays without PBC.

\begin{enumerate}
\item Sweep $N \in \{14, 20, 30, 40\}$ with and without PBC, at $Q \approx 1$ and at the two extremes;
\item Repeat domain and avalanche analyses; extrapolate the characteristic domain scale with $1/N$;
\item Quantify the systematic free-boundary vs.\ periodic deviation.
\end{enumerate}

\textbf{Destination:} appendix/methods section of the central article, or an independent methodological note. The $N = 40$ point without optimizations is expensive ($\mathcal{O}(K^2)$ with $K = 1600$), making this study the main motivation for the performance optimization proposed in Section~\ref{sec:melhorias}.

% ======================================================================
\section{Proposed code improvements}
\label{sec:melhorias}
% ======================================================================

Analysis of the V39 code identified the following opportunities, ordered by estimated impact. None blocks Study~1, which is already feasible; the performance items make Studies~2--4 substantially cheaper.

\subsection{Performance}

\begin{enumerate}
\item \textbf{Precomputation of the dipolar-interaction tensor --- IMPLEMENTED (V40, 04/07).} Needle positions are fixed and only the angles evolve, but up to V39 the $K \times K$ geometry matrices (displacements, distances, powers of $r$, \emph{cutoff} masks, and PBC image sums) were rebuilt at every integration step. Because the dipolar field is \emph{linear} in the moments $(m_x, m_y)$, all geometry was condensed once before the loop into three constant matrices $A_{xx}, A_{xy}, A_{yy}$ ($A_{yx} = A_{xy}$; periodic images already summed), reducing each step to matrix--vector products via BLAS/cuBLAS.

\emph{Validation} (three levels): tensor torques vs. stepwise method identical to machine \emph{epsilon} (maximum relative error $3 \times 10^{-16}$, for the 3 geometries, with and without PBC); full hysteresis trajectory with 172 steps and bit-identical $t$ and $B$ relative to V39, with $M$ and $S$ agreeing within $5 \times 10^{-12}$ (expected accumulation from floating-point reassociation); all 4 field modes and \mbox{Ctrl+C} behavior (code 130) preserved.

\emph{Measured result} (array $30\times30$, same baseline parameters): from 117.8~s / 148.0~s per hysteresis run to \textbf{3.4~s / 3.3~s}, a \textbf{35--45$\times$} speedup, above the initial 5--20$\times$ estimate. The tensor occupies only 19~MB for $K = 900$. Built-in memory protection: above an estimated 4~GB (K~$\gtrsim$~13,000), the code reports the issue and automatically falls back to the per-step method. The full campaign dropped from $\sim$7--8~h to $\sim$15~min serial;
\item \textbf{Campaign parallelization --- IMPLEMENTED (\code{damping\_sweep} V02, 04/07).} Three features: \code{--n\_workers~N} runs independent processes (each worker in its own working directory, avoiding collisions among incidental files written by the simulator); \textbf{incremental} \emph{manifest}, updated after each completed run (the previous version wrote it only at the end, so an interruption lost the whole record, a vulnerability that appeared during tests); and \code{--resume~1}, which resumes an interrupted campaign by skipping already registered runs.

\emph{Validation}: parallel execution with 4 \emph{workers} produces CSVs bit-identical to serial execution; full resume (``nothing to do'') and partial resume (simulated interruption) work; the analysis script reads the V02 \emph{manifest} without changes. The time gain is expected to be approximately linear in the number of cores, not measurable in the one-core development environment and to be confirmed on the server;
\item \textbf{Optional single precision on GPU (proposed).} \code{float32} doubles memory throughput; it should be sufficient for avalanche statistics, but must be validated against \code{float64} in one case before adoption;
\item \textbf{Adaptive time step in the overdamped regime (proposed; investigate cautiously).} With $Q \ll 1$, the dynamics is slow and $\Delta t$ tied to $T_0$ is conservative. A larger step, limited by the viscous scale, could reduce the cost of the slowest runs. This requires a dedicated convergence study before adoption. With the V40 gain, it is no longer a priority.
\end{enumerate}

Two relevant optimizations are already implemented and should not be redone: vectorized rendering (17--30$\times$, V35) and NVENC video encoding when \code{--gpu 1}.

\subsection{Visual appearance}

\begin{enumerate}
\item \textbf{Publication-figure mode} (\code{--fig\_style paper}): white background, colorblind-safe palette, vector export (PDF/SVG), fonts sized for article columns, and no HUD elements (clock, field indicator). Current figures (dark theme, PNG) are excellent for screen/video, but unsuitable for the manuscript. This is the highest immediate-value visual improvement for publication;
\item \textbf{Live $M\times H$ curve in the hysteresis video --- IMPLEMENTED (V45).} The \code{hysteresis} video now has a dual panel: needles on the left and the $M\times H$ curve on the right, traced in real time with the current point highlighted, communicating the full experiment in a single video. Other video modes remain single-panel; physics is bit-identical, since this is a rendering-only change;
\item \textbf{Visual highlighting of avalanches}: emphasize needles that rotated above a threshold in the last $n$ steps, showing the spatial propagation of avalanches and connecting the video to the paper statistics;
\item \textbf{Domain-wall drawing}: draw edges between domains with distinct labels, complementing colored halos and improving readability in large arrays;
\item \textbf{Net-magnetization arrow}: global vector $\vec M$ overlaid on the array, with instantaneous magnitude and direction;
\item \textbf{Energy halo mode}: heat map of dipolar energy per needle, revealing local frustration and likely avalanche nuclei.
\end{enumerate}

% ======================================================================
\section{Next steps}
% ======================================================================

\begin{enumerate}
\item Run the full campaign on the server (command documented in the scripts; \code{--quick\_test 1} is recommended first to validate the environment);
\item Run the analysis and apply the A/B criterion to $\tau(Q)$ with full statistics; in particular, check whether the slope of $\tau(Q)$ differs between geometries and whether the triangular-domain non-monotonicity survives;
\item Robustness diagnostics before any conclusion: convergence in $\Delta t$ at the extremes of $Q$; finite-size effects ($N = 20$ and $40$ at $Q \approx 1$); comparison with/without PBC;
\item Choose the journal based on the result (A $\to$ ACS Materials Au / Physical Chemistry Au; B $\to$ consider PRL) and start the manuscript framed as a ``model for \emph{artificial spin ice} and patterned nanomagnets'';
\item Satellite studies from the original roadmap, according to capacity: domain-formation kinetics by geometry; memory of the ``director'' after a pulse; effect of PBC on domain scale.
\end{enumerate}

% ======================================================================
\section{Version record}
% ======================================================================

\begin{longtable}{@{}lp{11.5cm}@{}}
\caption{Version record for \code{compass\_sim}. Versions V1--V19 precede the formal record and were consolidated retroactively; from V20 onward each change generates a new numbered version.}
\label{tab:versoes} \\
\toprule
\textbf{Version} & \textbf{Main content} \\
\midrule
\endfirsthead
\toprule
\textbf{Version} & \textbf{Main content} \\
\midrule
\endhead
\bottomrule
\endfoot
\bottomrule
\endlastfoot
V1--V19 & Simulator core: inertial dipolar dynamics (SI), three geometries, external field, MP4 video, field modes (static/pulse/hysteresis/sine), CSV export, PBC, and corrections to honeycomb and \code{t\_sim} semantics. \\
V20 & Formal version-record milestone. \\
V21--V24 & Pulse-mode refinement ($S$ plateau criterion); repositioned on-screen indicators (field at upper right, dedicated upper margin). \\
V25--V28 & GPU support via CuPy; \code{--gpu} flag; utilization investigation; performance instrumentation (throughput, ms/step). \\
V29--V33 & PBC with structure images in the dipolar field (default one copy); progress bar; automatic moment of inertia from real geometry (\SI{0.4}{mm} steel); Ctrl+C aborts immediately / Ctrl+I saves and exits (code 130). \\
V34 & Integration of GPU vectorized torque loop. \\
V35 & Matplotlib rendering optimization through vectorized collections: 17--30$\times$ speedup; bitwise validation of geometry and physics. \\
V36 & Parameter \code{--progress\_bar} \{0,1\}. \\
V37 & Automatic magnetic moment ($m = M_s V$, saturated steel); correction of status-message regression with \code{--progress\_bar 0}. \\
V38 & Magnetic-domain visualization: \code{--halo\_mode domains}, \code{--domain\_tol}; \emph{union--find} labeling; small domains in gray. \\
V39 & Parameter \code{--make\_images} \{0,1\}: suppresses all PNG generation while preserving CSVs; $\sim 3\times$ speedup in small arrays. \\
V40 & Precomputed dipolar-interaction tensor: matrix--vector products (BLAS/cuBLAS) per step; measured speedup of \textbf{35--45$\times$}; validation to machine \emph{epsilon}; memory protection with automatic fallback. \\
V41 & New field protocols: hysteresis with log-symmetric spacing (\code{--hyst\_spacing log}, $\sim$15$\times$ finer sampling near $B{=}0$ for $k{=}5$); pulse with explicit timing (\code{--field\_delay}, \code{--t\_pulse}) and torque-equilibrium criterion (\code{--torque\_tol}); step modes \code{step\_pos}/\code{step\_neg}. Legacy modes bit-identical to V40. \\
V42 & Parameter \code{--t\_sim\_full} \{0,1\}: with 1, the simulation runs until the end of \code{t\_sim}, disabling all physical early stops ($S{=}1.00$, rest, torque equilibrium), useful for uniform-length time series in sweeps. Ctrl+I and Ctrl+C remain active; default (0) bit-identical to V41. \\
V43 & User-reported fix: the comparison figure did not draw halos in any mode. Halo logic was extracted to \code{draw\_halos\_batch()} and applied to both comparison panels, with domain count in each panel title in \code{domains} mode. Also corrected the comparison supertitle, which displayed ``B\_ext=0.00'' for mT fields. \code{plot\_state} figures are pixel-identical and physics is bit-identical to V42. \\
V44 & Two-line live progress panel: bar and periodic status are rewritten in place using ANSI codes with a cursor invariant, eliminating terminal line stacking. TTY detection preserves log behavior without escape codes for redirected output. Validated with \code{pyte}; final screen with exactly one bar line and one status line; clean interleaved phase-transition messages; bit-identical physics. \\
V45 & Hysteresis-mode video with dual panel: needles on the left and the $M\times H$ curve traced in real time on the right, with the current point highlighted. Other video modes unchanged; bit-identical physics. \\
V46 & Alignment of help-text frames in the \emph{docstring}: right borders of parameter boxes, which meandered in plain-text editors because of visual-width differences (accented characters, arrows, \emph{box-drawing}), were normalized to uniform visual width; three long descriptions were slightly shortened to fit. No effect on physics or execution. \\
V47 & Autoscale in the $M\times H$ hysteresis curve: added parameter \code{--hyst\_autoscale} to dynamically scale the axes in both the live simulation frames (dual panel) and the final hysteresis plots, based on the accumulated minimum and maximum fields and magnetization. Translated all new comments and help texts to English. \\
V48 & Bugfix and parameter adjustment: resolved a syntax error on \code{box\_w} assignment in the screen indicators; changed the parameter \code{--hyst\_autoscale} to accept integers \code{0} and \code{1} (with default \code{0}) instead of a boolean flag. \\
V49 & Parameter and visualization refinement: resolved overlap between HUD boxes and the top row of needles by dynamically reserving the top 20\% of the vertical space; adjusted the field indicator bar to accept and represent signed values, ranging from $-B_{\max}$ (empty) to $+B_{\max}$ (full), with the center (half-full) representing zero field; updated the box arrow to rotate by $180^\circ$ for negative fields. \\
V50 & Progress bar translation and English layout: reorganized the terminal output progress tracking to a clean 3-line print format using standard English labels (`Integrating`, `step`, `steps/s`) to display runtime physics parameters and step counts in real time. \emph{(Current version.)} \\
\midrule
\multicolumn{2}{@{}l}{\emph{Satellite scripts (own version record):}} \\
\code{damping\_sweep} V01 & orchestration of the damping~$\times$~seed~$\times$~geometry campaign. \\
\code{damping\_sweep} V02 & parallel execution (\code{--n\_workers}), incremental \emph{manifest}, and resume (\code{--resume}). \\
\code{damping\_sweep\_analysis.py} & loop metrics, avalanche detection, MLE power-law fitting, plots. \\
\end{longtable}

\appendix
\section{Appendix: command-line parameters (V50)}
% ======================================================================

Complete list of parameters accepted by \code{compass\_sim} in version V50, extracted directly from the program interface. Parameters marked ``auto'' are calculated from the real physical geometry when omitted and overridden when explicitly supplied.

\begingroup
\small
\setlength{\tabcolsep}{5pt}
\renewcommand{\arraystretch}{1.15}

\begin{longtable}{@{}lll p{7.6cm}@{}}
\caption{Command-line parameters of \code{compass\_sim} V50.}
\label{tab:parametros}\\
\toprule
\textbf{Parameter} & \textbf{Default} & \textbf{Unit} &
\textbf{Description} \\
\midrule
\endfirsthead
\toprule
\textbf{Parameter} & \textbf{Default} & \textbf{Unit} &
\textbf{Description} \\
\midrule
\endhead
\bottomrule
\endlastfoot
\multicolumn{4}{@{}l}{\emph{Lattice geometry}} \\
\midrule
\code{--N} & 8 & --- & Number of needle rows. \\
\code{--M} & 8 & --- & Number of needle columns. \\
\code{--R} & 0.025 & m & Radius of the circle circumscribing each needle; the center-to-center distance between neighbors is $2R$ in all geometries. \\
\code{--needle\_frac} & 0.80 & --- & Needle length as a fraction of the diameter $2R$ (limited to 0--0.8). \\
\code{--geometry} & square & --- & Lattice type: \code{square}, \code{triangular}, or \code{honeycomb}. \\
\midrule
\multicolumn{4}{@{}l}{\emph{Physics and simulation}} \\
\midrule
\code{--moment} & auto & A$\cdot$m$^2$ & Magnetic moment per needle. If omitted: $m = M_s V$, steel saturated along the long axis (uses \code{--steel\_Bsat} and the real volume). \\
\code{--inertia} & auto & kg$\cdot$m$^2$ & Moment of inertia per needle. If omitted: real steel sheet, $I = \tfrac{1}{12}\,\mathrm{mass}\,(L^2 + \mathrm{width}^2)$. \\
\code{--needle\_thickness} & $4 \times 10^{-4}$ & m & Steel-sheet thickness used in both automatic calculations. \\
\code{--steel\_density} & 7850 & kg/m³ & Steel density (automatic inertia calculation). \\
\code{--steel\_Bsat} & 2.0 & T & Saturation flux density of steel (automatic moment calculation; $M_s = B_{\mathrm{sat}}/\mu_0$). Typical range: 1.6--2.2~T. \\
\code{--damping} & $5 \times 10^{-8}$ & N$\cdot$m$\cdot$s/rad & Viscous damping coefficient $b$. Zero = perpetual oscillation. \\
\code{--t\_sim} & 2.0 & s & Total simulated physical time (sum of integrated $\Delta t$ values). \\
\code{--t\_sim\_full} & 0 & --- & 1 = run until the end of \code{t\_sim}, disabling physical early stops ($S{=}1.00$, rest, torque equilibrium); Ctrl+I/Ctrl+C remain active. Useful for uniform-length sweep series. \\
\code{--field\_mode} & static & --- & \code{static}, \code{hysteresis} (five-ramp cycle), \code{sine}, \code{pulse} (wait $\to$ pulse $\to$ stabilization), \code{step\_pos} ($B{=}0$ for \code{field\_delay}, then kept on), or \code{step\_neg} (field from $t{=}0$, removed at \code{field\_delay}). \\
\code{--hyst\_spacing} & linear & --- & Hysteresis-ramp spacing: \code{linear} or \code{log} (log-symmetric through sinh, with small $dB/dt$ near $B{=}0$, fine sampling in the coercivity region, large near $\pm B_{\max}$). \\
\code{--hyst\_log\_k} & 5.0 & --- & Concentration of \code{log} spacing: $k \to 0$ recovers linear; $k{=}5$ concentrates $\sim$15$\times$ more points near $B{=}0$. \\
\code{--field\_delay} & 0.0 & s & Waiting time $\Delta t$: in \code{pulse}, time with $B{=}0$ before the pulse; in \code{step\_pos}, before switching on; in \code{step\_neg}, time with field on before removal. \\
\code{--t\_pulse} & --- & s & Pulse duration (\code{pulse}). If omitted, legacy criterion: field until $S \geq 0.99$ or $S$ plateau. \\
\code{--torque\_tol} & $10^{-3}$ & --- & Relative tolerance of equilibrium criterion in pulse/step modes: stop when $\omega_{\max}$ is small \emph{and} $\langle|\tau|\rangle < \mathrm{tol}\cdot\tau_{\mathrm{ref}}$ ($\tau_{\mathrm{ref}} = I\omega_0^2$) for 50 steps; 0 disables torque criterion. \\
\code{--field\_freq} & 1.0 & Hz & Sinusoidal-field frequency (only \code{sine}). \\
\code{--dt\_factor} & 0.05 & --- & Fraction of natural period $T_0$ used as time step $\Delta t$ (0.02--0.10). \\
\code{--noise} & 1.5 & rad & Initial angular noise amplitude ($0$ = all along $+x$; $\pi$ = fully random). \\
\code{--seed} & 42 & --- & Random-number-generator seed (reproducibility). \\
\code{--pbc} & 0 & --- & Periodic boundary conditions (0/1); with 1, edges interact with replicas through the minimum-image convention. \\
\code{--pbc\_images} & 1 & --- & Replicas summed on each side in each direction when \code{--pbc 1} (default: $3\times3$ cell grid). \\
\code{--gpu} & 0 & --- & GPU backend via CuPy (0/1), with automatic fallback to CPU if unavailable. \\
\code{--progress\_bar} & 1 & --- & Real-time progress bar (0/1); turn off in redirected logs. \\
\code{--halo\_mode} & order & --- & Halo coloring: \code{order} (local alignment) or \code{domains} (connected domains, V38). \\
\code{--domain\_tol} & 15.0 & degrees & Angular tolerance for neighbors to belong to the same domain (\code{--halo\_mode domains}). \\
\midrule
\multicolumn{4}{@{}l}{\emph{Uniform external field}} \\
\midrule
\code{--B\_ext} & 0.0 & T & External-field magnitude (e.g., $50 \times 10^{-6}$ = Earth's field). \\
\code{--phi\_ext} & 0.0 & degrees & Field direction (0 = $+x$, 90 = $+y$, counterclockwise). \\
\code{--ext\_Bx}, \code{--ext\_By} & --- & T & Explicit components; override \code{--B\_ext}/\code{--phi\_ext}. \\
\midrule
\multicolumn{4}{@{}l}{\emph{Output}} \\
\midrule
\code{--video} & --- & --- & Base name of the MP4 video (ffmpeg); automatic numbering \code{name0000.mp4}, \code{0001}, \dots \\
\code{--frame\_every} & 5 & steps & Interval between saved frames. \\
\code{--make\_images} & 1 & --- & 0 = disables \emph{all} PNG generation (frames and summaries), while keeping data CSVs (V39). \\
\code{--fps} & 24 & --- & Video frames per second. \\
\code{--dpi} & 120 & --- & Frame resolution (arrays $>15\times15$: use 150--200). \\
\code{--keep\_frames} & off & --- & Keeps the intermediate PNG folder after MP4 generation. \\
\code{--hyst\_autoscale} & 0 & --- & Enables automatic scaling of the axes in both the live simulation frames (dual panel) and the final hysteresis plots (0/1; V50). \\
\code{--csv\_order} & t & --- & CSV column order: \code{t} $\to$ $t,B,M,S$; \code{B} $\to$ $B,M,S,t$. \\
\end{longtable}
\endgroup

% ======================================================================
\section{Appendix: usage examples}
\label{app:exemplos}
% ======================================================================

Command-line recipes for the most frequent use cases, checked against version V50. Commands use \code{compass.py} as the executable name, a copy of the current simulator version (\code{compass\_simV50.py}) kept under this stable name for daily use, so the examples can be copied and pasted directly. A final backslash indicates shell line continuation.

\subsection*{B.1 --- Basic execution (dipolar equilibrium)}

Square $8\times8$ array with all defaults (steel needle with automatically computed moment and inertia, no external field):
\begin{verbatim}
python3 compass.py --N 8 --M 8 --geometry square
\end{verbatim}
Generates summary images (initial state, equilibrium, comparison, order parameter) and \code{compass\_field\_log.csv} in the current directory. For the other lattices: \code{--geometry triangular} or \code{--geometry honeycomb}.

\subsection*{B.2 --- GPU on and off}

CPU is the default (\code{--gpu 0}). With V40 optimization, the CPU already handles medium arrays in seconds; GPU is useful for large arrays:
\begin{verbatim}
# CPU (default) -- medium array
python3 compass.py --N 30 --M 30 --t_sim 5.0

# GPU (CuPy/RTX 3090) -- large array; requires cupy installed
python3 compass.py --N 60 --M 60 --t_sim 5.0 --gpu 1
\end{verbatim}
With \code{--gpu 1}, video encoding also uses the NVENC hardware encoder when available. If CuPy is not functional, the program warns and automatically falls back to CPU.

\subsection*{B.3 --- Video and figures in good resolution}

For arrays above $15\times15$, use \code{--dpi} 150--200. Smaller \code{--frame\_every} = smoother animation (more frames):
\begin{verbatim}
python3 compass.py --N 20 --M 20 --t_sim 4.0 \
    --video demo --dpi 180 --fps 30 --frame_every 3
\end{verbatim}
The video receives automatic numbering (\code{demo0000.mp4}, \code{demo0001.mp4}, \dots). To keep individual frame PNGs after assembling the video, add \code{--keep\_frames}. For static high-resolution figures only (no video), omit \code{--video} and use \code{--dpi 200}.

\subsection*{B.4 --- Fast sweeps: no images, no bar}

In campaigns with many executions, turn off all image generation (data CSVs are kept) and the progress bar:
\begin{verbatim}
python3 compass.py --N 8 --M 8 --t_sim 1.0 \
    --make_images 0 --progress_bar 0
\end{verbatim}
Speedup of $\sim 3\times$ in small arrays, where rendering dominates time. For \emph{uniform}-length time series across runs (no early stops by $S{=}1.00$, rest, or equilibrium), add \code{--t\_sim\_full 1} (V42).

\subsection*{B.5 --- Damping control (dynamic regimes)}

The regime is characterized by the quality factor $Q$, printed in the header of each run along with the classification (under-/overdamped). The exact $Q$ depends on the dominant field (dipolar or external); the values below, with all other parameters at default, illustrate the three regimes:
\begin{verbatim}
# underdamped (Q >> 1): visibly oscillates before stabilizing
python3 compass.py --N 10 --M 10 --damping 5e-8 --video sub

# intermediate (Q ~ 1): crossover
python3 compass.py --N 10 --M 10 --damping 1e-6 --video middle

# overdamped (Q << 1): monotonic relaxation, no oscillation
python3 compass.py --N 10 --M 10 --damping 5e-5 --video over
\end{verbatim}
Check the printed $Q$ in the header and adjust \code{--damping} until the desired regime is reached. \code{--damping 0} disables dissipation (perpetual oscillation).

\subsection*{B.6 --- Needle size and spacing}

The radius $R$ of the circumscribed circle defines lattice spacing ($r_{nn} = 2R$); \code{--needle\_frac} defines needle length as a fraction of $2R$; sheet thickness affects automatic inertia and moment calculations:
\begin{verbatim}
# needles with 1 cm radius, occupying half the available diameter
python3 compass.py --N 12 --M 12 --R 0.01 --needle_frac 0.5

# thinner sheet (0.2 mm): lighter needle -> oscillates faster
python3 compass.py --N 12 --M 12 --needle_thickness 2e-4

# override automatic calculations with explicit values
python3 compass.py --N 12 --M 12 --moment 0.05 --inertia 1e-9
\end{verbatim}

\subsection*{B.7 --- External field: the six modes}

One example per mode, covering the full range of \code{--field\_mode}. The amplitude is \code{--B\_ext} and the direction is \code{--phi\_ext}.

\medskip\noindent\textbf{(1) static} --- constant (or zero) field throughout the simulation:
\begin{verbatim}
python3 compass.py --N 10 --M 10 --B_ext 1e-3 --phi_ext 45
\end{verbatim}

\noindent\textbf{(2) hysteresis} --- five-ramp cycle ($0 \to +B_{\max} \to 0 \to -B_{\max} \to 0 \to +B_{\max}$). The video (V45) shows needles and the $M\times H$ curve traced in real time, side by side. Spacing can be \code{linear} (default) or \code{log} ($\sim$15$\times$ finer sampling in the coercivity region):
\begin{verbatim}
# linear, with dual-panel video (needles + M x H curve)
python3 compass.py --N 20 --M 20 --field_mode hysteresis \
    --B_ext 0.002 --t_sim 10.0 --video hist --dpi 160

# log-symmetric spacing (fine near B=0)
python3 compass.py --N 20 --M 20 --field_mode hysteresis \
    --hyst_spacing log --hyst_log_k 5 --B_ext 0.002 --t_sim 10.0
\end{verbatim}

\noindent\textbf{(3) sine} --- sinusoidal field $B(t) = B_{\max}\sin(2\pi f t)$; \code{--t\_sim} should cover several periods:
\begin{verbatim}
python3 compass.py --N 15 --M 15 --field_mode sine \
    --B_ext 0.002 --field_freq 0.5 --t_sim 8.0
\end{verbatim}

\noindent\textbf{(4) pulse} --- wait \code{field\_delay}, apply for \code{t\_pulse}, set to zero, and stabilize (equilibrium by $\omega$ and torque). If \code{--t\_pulse} is omitted, the legacy criterion is used (field until $S \geq 0.99$ or plateau):
\begin{verbatim}
# timed: wait 0.2 s -> field 0.3 s -> zero and stabilize
python3 compass.py --N 15 --M 15 --field_mode pulse \
    --field_delay 0.2 --t_pulse 0.3 --B_ext 0.005 --t_sim 6.0

# legacy (without t_pulse): field until saturation, then relaxation
python3 compass.py --N 15 --M 15 --field_mode pulse \
    --B_ext 0.005 --t_sim 6.0
\end{verbatim}

\noindent\textbf{(5) step\_pos} --- $B{=}0$ for \code{field\_delay}, then field switched on and held until the end:
\begin{verbatim}
python3 compass.py --N 15 --M 15 --field_mode step_pos \
    --field_delay 0.3 --B_ext 0.003 --t_sim 3.0
\end{verbatim}

\noindent\textbf{(6) step\_neg} --- field from $t{=}0$, removed at \code{field\_delay}, then kept at zero until the end:
\begin{verbatim}
python3 compass.py --N 15 --M 15 --field_mode step_neg \
    --field_delay 0.3 --B_ext 0.003 --t_sim 3.0
\end{verbatim}

\noindent In \code{hysteresis}, the $M(B)$ loop is exported to \code{hysteresis\_loop.csv}; in \code{sine}, to \code{sine\_field.csv}. In pulse/step modes, \code{--torque\_tol} controls the equilibrium criterion and \code{--t\_sim\_full 1} disables early stopping (Appendix~A).

\subsection*{B.8 --- Magnetic-domain visualization}

Demonstration scenario producing clear, well-separated domains (square lattice with low initial noise and long relaxation):
\begin{verbatim}
python3 compass.py --N 24 --M 24 --R 0.01 --noise 0.6 \
    --t_sim 8.0 --damping 3e-6 \
    --halo_mode domains --domain_tol 35 --video domains --dpi 160
\end{verbatim}
Each connected region of needles with similar orientation (within \code{--domain\_tol} degrees) receives a color; small domains appear in gray. Smaller tolerances, for example the default $15^\circ$, fragment more; larger tolerances aggregate more. From V43 onward, halos (domain or order) appear in \emph{all} figures: \code{compass\_initial}, \code{compass\_equilibrium}, video frames, and also \code{compass\_comparison}, whose panels show the domain count for each state in the title.

\subsection*{B.9 --- Periodic boundary conditions and reproducibility}

\begin{verbatim}
# array treated as an infinite periodic structure (1 replica per side)
python3 compass.py --N 16 --M 16 --pbc 1 --pbc_images 1

# exact reproducibility: fix the initial-noise seed
python3 compass.py --N 10 --M 10 --seed 123
\end{verbatim}

\subsection*{B.10 --- Full scientific campaign (satellite scripts)}

\begin{verbatim}
# quick environment validation (seconds)
python3 damping_sweepV02.py --out_dir results --quick_test 1

# full Study 1 campaign, parallel on 8 cores
python3 damping_sweepV02.py --out_dir results \
    --grid_n 30 --n_seeds 5 --n_dampings 8 \
    --Q_min 0.05 --Q_max 15.0 \
    --geometries square,triangular,honeycomb --n_workers 8

# resume interrupted campaign (skip already registered runs)
python3 damping_sweepV02.py --out_dir results ... --resume 1

# analysis: loop metrics, avalanches, tau(Q), domains
python3 damping_sweep_analysis.py --sweep_dir results \
    --out_dir analysis
\end{verbatim}
Avoid combining \code{--n\_workers} $> 1$ with \code{--use\_gpu 1}, due to GPU contention between processes.

% ======================================================================
\section*{Version history of this report}
\addcontentsline{toc}{section}{Version history of this report}
% ======================================================================

\begin{tabular}{@{}lll@{}}
\toprule
\textbf{Version} & \textbf{Date} & \textbf{Changes} \\
\midrule
V01 & 04/07/2026 & Initial version. \\
V02 & 04/07/2026 & Added Appendix~A (command-line parameters of V39). \\
V03 & 04/07/2026 & Added Sections~\ref{sec:roteiro} (step-by-step roadmap of proposed studies) and \ref{sec:melhorias} (proposed code improvements: performance and visual). \\
V04 & 04/07/2026 & Documented implementation of performance improvements 1 and 2 (V40 code: precomputed tensor, 35--45$\times$; \code{damping\_sweep} V02: parallelization, incremental \emph{manifest}, resume), with validations and measurements; updated campaign costs and version record. \\
V05 & 04/07/2026 & Added Appendix~B (usage examples: GPU, high-resolution video/figures, damping regimes, needle geometry, field modes, domains, PBC, and full campaign). \\
V06 & 04/07/2026 & Documented code V41: hysteresis with log-symmetric spacing, pulse with explicit timing ($\Delta t$ + $t_{\mathrm{pulse}}$ + torque equilibrium), and step modes \code{step\_pos}/\code{step\_neg}; updated Section~2, version record, and Appendices A and B. \\
V07 & 04/07/2026 & Documented code V42: parameter \code{--t\_sim\_full} (run until the end of \code{t\_sim}, without physical early stops). \\
V08 & 04/07/2026 & Documented code V43: halos (order and domains) now also in the comparison figure, with domain count by panel; correction of \code{B\_ext} formatting in the comparison supertitle. \\
V09 & 04/07/2026 & Documented code V44: two-line progress panel (bar + status) updated in place, with TTY detection. \\
V10 & 05/07/2026 & Documented V45 (hysteresis video with dual panel: needles + $M\times H$). Editorial revision: Appendix~B.7 rewritten as reference for the six field modes without redundancy; example-command versions standardized to V45; Appendix~A title updated; Phase~6 chronology extended to V45. \\
V11 & 05/07/2026 & In Appendix~B examples, the executable is now referred to as \code{compass.py} (stable copy of the current version), for direct copy and paste. The version record (V01--V45) is preserved in the body and tables. \\
V12 & 05/07/2026 & Documented V46 (alignment of help-text frames). \\
V13 & 05/07/2026 & Moved the ``Version record'' section before the appendix and converted it to a \code{longtable} to allow multi-page spanning; documented code V47 (hysteresis autoscale parameter \code{--hyst\_autoscale} and English comments translation). \\
V14 & 05/07/2026 & Documented code V48 (syntax bugfix in \code{box\_w} and adjustment of \code{--hyst\_autoscale} to integer choices 0/1). \\
V15 & 05/07/2026 & Documented code V49 (HUD overlap fixes, bipolar signed field bar, and arrow rotation updates). \\
V16 & 05/07/2026 & Documented code V50 (re-styled English progress bar on a 3-line terminal console layout). \\
\bottomrule
\end{tabular}

% ======================================================================
\begin{thebibliography}{9}

\bibitem{dahmen2011}
K.~A. Dahmen, Y.~Ben-Zion, J.~T. Uhl,
\emph{A simple analytic theory for the statistics of avalanches in
sheared granular materials},
Nature Physics \textbf{7}, 554 (2011).

\bibitem{nisoli2013}
C.~Nisoli, R.~Moessner, P.~Schiffer,
\emph{Colloquium: Artificial spin ice: Designing and imaging magnetic
frustration},
Reviews of Modern Physics \textbf{85}, 1473 (2013).

\bibitem{prb1998}
Square lattice of interacting magnetized needles: thermodynamic
properties through random field and Monte Carlo,
Physical Review B \textbf{58}, 9238 (1998).

\bibitem{clauset2009}
A.~Clauset, C.~R. Shalizi, M.~E.~J. Newman,
\emph{Power-law distributions in empirical data},
SIAM Review \textbf{51}, 661 (2009).

\end{thebibliography}

\end{document}
